\reviewsection{Whole-Person and Multiple-Identity Lens}

As I argued in my Module 1 journal, developmental difference is neurological, but it is also lived, relational, cultural, historical, and socially interpreted. My own experience as a Latino, bilingual, neurodivergent person living with chronic illness and shaped by trauma has taught me that institutions often notice one visible difficulty before they notice context, effort, intelligence, language, culture, or persistence. That experience makes me cautious about allowing autism, ADHD, disability, behavior, or any other single category to become a complete explanation of a student.

Students also enter classrooms through multiple identities. An autistic learner may also be multilingual, racialized, gender-diverse, economically marginalized, physically disabled, chronically ill, or part of a family and culture with its own language for disability and support. Those identities affect who is noticed, believed, disciplined, evaluated, included, or given access to services. A peer lesson about autism therefore cannot teach students to replace one stereotype with a more positive stereotype. It must teach curiosity, consent, privacy, and respect for variation across identities.

This is where the lesson becomes part of my Puzzle Plan. No diagnosis, behavior, peer observation, family story, cultural interpretation, or classroom data point stands alone. Each is one piece. Through the EQUITAS lens, I ask whose knowledge is being centered, what context is missing, what power shapes the interpretation, and whether the response expands communication, dignity, agency, participation, and belonging. The goal is not to make a student easier for the classroom, but to make the classroom more responsive to the whole student.
